% PDF 343--349；原书第 316--322 页。
\section{大语言模型的应用}
\label{sec:ch08-applications}
讨论开发可用于生产环境的大语言模型应用的技术细节之前，先了解这类模型已用于哪些问题和场景。这里简要介绍大语言模型能够改善结果的几种通用应用，例如对话式聊天机器人，再概述可用于这些场景的不同类型的模型。

这些应用可以帮助读者把本章较技术性的内容放回开发周期中理解，看到它们在实践中的用途。

\subsection{对话式人工智能、聊天机器人和人工智能助手}
这是大语言模型带来显著改进的最常见场景。在聊天机器人和对话式人工智能中，与传统语言模型相比，大语言模型有更强的语言理解能力，带来了几项重要的新能力，其中尤为重要的是自然语言理解（natural language understanding，NLU）（\refcite{ch08-ref-098}{Wei et al., 2022}）。在这类应用中，大语言模型能否理解用户查询背后的意图，并将输入与已有参数知识综合起来，生成连贯的回答，对应用能否发挥作用影响很大。

用户通常进行开放式对话，可能涉及多个知识领域或主题，因此大语言模型能否跟踪上下文，也关系到整个对话会话中的回答是否连贯。在用户与应用来回交流的多轮对话中，大语言模型需要有选择地纳入前面的查询或回答，从而始终给出有用、连贯的回答。输入上下文长度的提升进一步增强了这一能力（\refcite{ch08-ref-069}{Pawar et al., 2024}），因为模型可以在更长的输入范围上进行推理。

\begin{keypoint}
这类大语言模型应用中的许多场景，对回答的事实性有较高要求。开发时需要尽量减少模型容易产生的幻觉。常见办法是使用外部知识库，从中提取相关上下文，以经过核查的知识作为模型回答的条件。怎样把知识库纳入应用架构，并把其中的知识放入模型输入以得到合适的回答，也给应用开发带来一组问题。
\end{keypoint}

\subsection{内容创作}
在内容营销和数字媒体领域，有一句常说的话叫“内容为王”。以往，故事、博客文章、通讯、社交媒体内容等，主要由有经验的人创作。他们知道哪些内容能引起受众共鸣，制作这些内容，并通过合适的渠道传播。当时已有不少内容制作环节可由大语言模型完成。有些应用允许营销人员整理有关受众人口特征的背景资料，提供相关内容和详细指导，由含有大语言模型的系统生成有吸引力的内容，并按预定日程在各渠道发布。

同样，可以用提示让大语言模型生成整篇文章，或创作关于真实或虚构主题与事件的故事。读者常常难以将这些内容与人创作的内容区分开来，内容创作者可以借此扩大创作规模，增加内容的多样性。这些发展也带来负面后果，在教育中尤其明显。学生很快开始用 OpenAI 的 ChatGPT 等大语言模型完成作业，可能未能充分掌握知识（\refcite{ch08-ref-053}{Lo, 2023}）。另一方面，大语言模型大幅提高了教育工作者设计、规划和制作课程内容的效率。使用得当时，它也是方便的学习辅助工具。

\subsection{搜索、信息检索和推荐系统}
这类应用主要是尽可能高效地为用户提供相关信息。检索常常像大海捞针。如果没有识别并呈现相关内容的办法，用户就会不断收到不切题或无关的信息。以一家公司为例，它有 10 年的产品文档，涉及多个产品版本及其功能，在这样的知识库中查找信息会很复杂。借助大语言模型，用户可以与人工智能助手对话。助手使用整合用户查询、知识库及其元数据的技术，使用户更高效地了解产品功能。

大语言模型的自然语言理解、实体抽取、摘要和语义嵌入能力有所改善，因此能以更好的精确率和召回率找出相关知识，进行推理，并将结果纳入个性化回答。这些变化有赖于大语言模型在方法、架构和工具上的许多改进，例如第~\ref{ch:rag}~章介绍的检索增强生成技术。

\urldef\githuturl\url{https://madnight.github.io/githut/#/}
\subsection{编程}
大语言模型不仅善于生成自然语言，也能生成计算机编程语言。原书以当时常用的 GitHub Copilot 为例。它为人类程序员开发软件提供代码辅助，下面用其主要功能说明此类大语言模型应用可以提供的帮助。

\begin{itemize}
\item \textbf{代码自动补全。}可以简单到像传统 Tab 补全那样补全函数、方法和变量名，也可以根据对当前脚本已有代码的实时分析，建议完整的代码块。
\item \textbf{支持多种编程语言。}使开发人员能更高效地在多种语言间协同工作。这在全栈或专业领域应用开发中尤其有用，因为方案的不同组件可能用不同语言编写。例如，一名全栈开发人员用 JavaScript 为用户界面编写数据处理程序时，Copilot 可以用 Python 建议后端 API 的代码块，为前端提供数据。据原书当时的描述，GitHub Copilot 支持公开 GitHub 仓库中可用的编程语言，其在不同语言上的能力取决于这些语言在公开仓库代码中的比重。GitHut\footnote{\githuturl} 展示了不同编程语言在 GitHub 上的相对比重。
\item \textbf{自然语言理解。}用户可以用自然语言说明希望代码实现哪些功能，往往只需编写描述下一段代码用途的注释。如果当前脚本启用了 GitHub Copilot，它会推荐实现注释所述功能的代码。对许多问题，这类建议可以给出简洁的解法，并能利用正在开发的脚本的其他上下文，尤其在注释和文档写得较好的情况下。这个功能能提高效率。与其他大语言模型生成内容一样，用户仍需核查并测试推荐的代码，防止因模型生成错误而出问题。
\item \textbf{代码重构与调试。}GitHub Copilot 的另一个用途是提高重构和调试的效率。训练使用了大量代码，模型学到了同一问题或相似问题的多种解法，因此可以提供其他模式供用户考虑。对于定义 Python 类等重复性工作，Copilot 可以生成样板代码，使开发者把时间用于具体功能。在调试时，GitHub Copilot 可以结合代码和开发人员对预期功能的描述，解读执行错误或意外输出，帮助定位可能出问题的代码。它也可以用自然语言说明执行流程和代码功能，协助开发人员查找执行问题的成因。
\end{itemize}

这些能力有助于提高软件开发效率。以往，要写出可以完成预期功能的代码，可能需要查阅参考书，多次访问 Stack Overflow 或 GitHub Gists 等网站，并请同事审查代码。现在，编程助手让开发人员可以通过同一个界面获得类似的学习和反馈。编程助手也已经集成进 Visual Studio Code、Vim 和 JetBrains 等常用集成开发环境。%
\footnote{译注：原书将 Visual Studio Code 写成“Visual Basic Code”。}

\subsection{大语言模型的类别}
\label{sec:ch08-categories}
大语言模型的数量迅速增加，开源领域尤其如此（\refcite{ch08-ref-027}{Gao and Gao, 2023}）。训练数据、微调方法、提示、算力规模、学习目标和其他影响开发的因素都可以调整，由此产生许多改进，推动模型多种能力的提升。单靠一章很难综述这些变化。不过，有许多在线资源可以帮助读者辨别具体模型在任务与基准上的表现。例如，Hugging Face 的 \texttt{open\_llm\_leaderboard} 随新模型开发与发布汇总大语言模型的任务表现。%
\footnote{\url{https://huggingface.co/spaces/HuggingFaceH4/open_llm_leaderboard}}
这类工具常允许用户按关心的因素筛选，例如只比较采用某种架构或某种权重精度的大语言模型。

比较各个模型是为具体场景选型的重要步骤，了解大语言模型的主要类别也很重要。不同类别擅长的应用不同，例如对话、编程与检索。下面介绍几种主要类别，帮助读者作出选择。

\subsubsection{通用大语言模型}
通用大语言模型通常在来自互联网的大型语料库上训练。一些语料经过筛选，但整体上，这类训练集广泛涵盖互联网上的不同主题和领域。因此，模型擅长通用语言任务，例如自然语言理解、实体抽取、思维链与问答，但未必擅长特定领域的任务，例如金融数值推理。这类大语言模型用途最广，可以作为许多应用的起点。若模型在某些方面能力不足，不少模型还可以针对任务做专门化处理，不过开源模型通常比闭源模型更便于这样做。概括地说，可以从文本生成、自然语言理解、问答、翻译、分类、摘要和对话等方面考察这类模型。对多数用例，它们直接使用时的表现可能已足够。

\subsubsection{多模态大语言模型}
多模态大语言模型在不止一种“模态”（modality）的数据上训练。常见模态包括文本、音频、视频和图像（\refcite{ch08-ref-105}{Yin et al., 2024}）。在这些模态上训练模型，带来一组新的跨模态应用，正在成为生成式人工智能的新前沿（详见第~9~章）。第~\ref{sec:ch08-applications}~节讨论的大语言模型应用，还可以借助多模态大语言模型扩展，例如依据自然语言描述检索图像，或根据自然语言指令生成音频。将文本与图像生成结合，用来讲故事或编写产品说明等，可以提供更丰富、更具表现力的用户体验，正在受到模型开发者关注。总体而言，与传统大语言模型相比，这些模型规模很大，给采用和集成带来新的困难。研究者正在继续改善这类模型的任务表现和计算效率。

\subsubsection{多语言大语言模型}
多语言大语言模型在不止一种自然语言的文本数据上训练。它们已受到广泛研究，可用于翻译、多语言推理、多语言内容生成等任务。有些模型支持的语言种类很多。例如，BigScience 开发的开源模型 BLOOM（\refcite{ch08-ref-103}{Workshop et al., 2023}），按原书此处的描述，具有 1960 亿参数，在 46 种自然语言和 13 种编程语言的文本上训练。%
\footnote{译注：本书引用的 BigScience 论文题名与第~4.4.1.1~节所用的 BLOOM-176B 均指 1760 亿参数模型。}
这类模型适合跨语言任务，例如对多种语言的文本作推理（如 \refcite{ch08-ref-076}{Ranaldi et al., 2024}）。在应用中，多语言大语言模型已用于客服等沟通场景。这些场景常出现语码转换（code-switching），即在交际中交替使用不同自然语言（\refcite{ch08-ref-106}{Yong et al., 2023}）。

\subsubsection{领域专用大语言模型}
领域专用大语言模型在来自特定行业、领域或专业方向的数据上训练，数据经过有针对性的选择。训练的目的通常是让模型的能力适应领域的特点。当领域有许多难以按通用语境理解的行话，或者其内容分布与通用语境差异较大时，这一点尤其重要。例如，在生物医学领域，研究者开发了领域专用大语言模型 BioMistral（\refcite{ch08-ref-040}{Labrak et al., 2024}）。他们在网上大型生物医学研究文献库之一 PubMed Central 上，对 Mistral 模型做适应性预训练。由 70 亿参数的通用 Mistral 适配得到的 BioMistral 模型，在 10 项生物医学任务中的 9 项优于通用模型。教育、法律、经济、政治、科学和金融等领域也有专用大语言模型，可以作为开发许多领域应用的起点。
